\cotefile{157-3}
\scanpage{267}
\mspage{217}
$\bullet$ pour \unsure{tte|tt} $S$ dans $M$ une sous-catégorie $\unsure{F|\mathcal F}_{0S}$ de $\underline{\mathrm{Fib}}_{W}\unsure{\backslash S|/S|S}$ (je pense par exemple : $\underline{\unsure{\mathrm{Parf}|\mathrm{Perf}}}_{W}S$, voire $\underline{\mathrm{Parf}}_{W_{\unsure{w|\infty}}}S$, dans le cas où $M=\mathrm{Cat}$), telle que pour \unsure{tte} s-catégorie pleine \inserted{$F$} de $M_{/S}$ telle que
\[
F_{0S}\subset F\subset M_{/S}
\]
(a fortiori, pour une $F$ pleine telle que $\underline{\mathrm{Fib}}_{W}S\subset F\subset M_{/S}$), les inclusions \deleted{\unsure{définissent}} \inserted{induisent} des équivalences de catégories
\[
F_{0S}(W_{\unsure{S}})^{-1}\xrightarrow{\approx}
F(W_S)^{-1}\xrightarrow{\approx}
\begin{array}{c}
(M_{/S})W_{\unsure{S|J}}^{-1}\\[-2pt]
\Vert\scriptstyle{\text{déf}}\\[-2pt]
\mathrm{Hot}^{!}_{W}(S)
\end{array}.
\]

Le choix de $F_{0S}$ ne semble pas pouvoir se fixer de façon ``\unsure{la meilleure|le meilleur} possible'', une fois pour toutes dans le contexte axiomatique adopté ici. C'est un élément de structure supplémentaire sur $M$ (en plus du localiseur $W$), au $=$ titre que la donnée des \unsure{classes} de Q-fibrations et Q-cofibrations à la Quillen. J'ignore si les axiomes a) b) ci-dessus (Axiomes de la fibre homotopique) impliquent déjà
% archive sheet number: 267
\scanpage{268}
\mspage{218}
la variante de \textcircled{a} :

\textcircled{c} Pour \unsure{tte} sous-catégorie pleine $\unsure{F|\mathcal F}$ de $M_{/S}$ telle que
\[
\underline{\mathrm{Fib}}_{W}S\subset F\subset M_{/S}
\]
les inclusions induisent des équivalences
\[
\begin{array}{c}
\underline{\rule[-.24ex]{.08ex}{.48ex}\vphantom{F}(\underline{\mathrm{Fib}}_W S)W_S^{-1}\rule[-.24ex]{.08ex}{.48ex}}\\
\Vert\\
\mathrm{Ho}^{\unsure{\mathrm{lc}|\mathrm{li}}}_{W}(S)\\
\text{\deleted{\unsure{moyennant les}}}\\
\text{moyennant a)}
\end{array}
\xrightarrow{\approx} F(W_S)^{-1}
\xrightarrow{\approx}
\begin{array}{c}(M_{/S})(W_{S})^{-1}\\ \Vert\\ \mathrm{Ho}^{!}_{W}(S)\end{array}
\]

Mais je présume que pratiquement, quand on pourra vérifier \textcircled{a}, on pourra vérifier \textcircled{\unsure{b|c}} également.
% slanted
\margin{cf page 240 \unsure{ff|f}}
% slanted
\margin{Il faudrait ici développer un peu plus le formalisme de $\mathrm{Ho}_{W}^{\mathrm{lc}}(S)$ et des fibres homotopiques, sous l'hypothèse \textcircled{a} \textcircled{b} plus haut.}

5. \uline{Axiomatique des $W$-types d'homotopie relatifs non loc. constants} $\mathrm{Ho}^{\unsure{\mathrm{lc}|\mathrm{c}}}_{W}(S)$

On n'a rien dit \unsure{jusqu'en|jusque} ce point d'éventuels \inserted{$W$-}Types d'homotopie relatifs (sur un \unsure{objet} $S$ de $M$) qui ne seraient \uline{pas} ``$\mathrm{loc}^{\mathrm{t}}$ constants''. Il en était question \unsure{pourtant} dans loc. cit (VI, théorème \unsure{salade} 2, partie II du théorème\unsure{)}, pages \unsure{107|109}-113.). Je profite de l'occasion pour revenir là dessus. Il faut introduire, dans chaque $M_{/S}$, un deuxième localiseur
\[
W_S^0\subset\mathrm{Fl}(M_{/S}),\quad W_S^0\subset W_S
\]
(plus fin que $W_S$), ce qui permettra de définir
% archive sheet number: 268
\scanpage{269}
\mspage{219}
\[
\mathrm{Ho}_W(S)\overset{\text{déf}}{=}(M_{/S})(W_S^0)^{-1}.
\]
Qu'il n'y ait pas de choix \uline{canonique} qui s'impose ici, déduit de la seule structure $(M,W)$, est clair sur l'exemple de $(\mathrm{Cat},W)$, auquel cas on trouve \uline{deux} choix symétriques transformés l'un en l'autre par l'involution
\[
X\longrightarrow X^{\unsure{\circ|0}}\qquad(\mathrm{Cat},W)\xrightarrow{\unsure{\sim}}(\mathrm{Cat},W),
\]
à savoir $W_S^g$ et $W_S^d$. Le choix de l'un ou de l'autre, disons (pour fixer les idées) de $W_S^g$, est associé à une notion d'\uline{iso local}\unsure{``|''} (ou de ``\unsure{coïso|coiso|coïsom.} local\unsure{''|.} dans le cas de $W_S^d$) que je voudrais axiomatiser. Il semble bien que pour les principales catégories de modèles pour $\mathrm{Hot}$ (et pour ses variantes), on dispose de façon assez naturelle d'une telle notion d'iso local (sans l'embarras, le plus souvent, d'un \unsure{autom.|autome} involutif fournissant une notion symétrique \ldots)

On supposera donc, comme dans loc. cit. p. 94-\unsure{95|96}, donné une classe
\[
\unsure{\mathcal L|L}\subset\mathrm{Fl}(M)
\]
de morphismes, jouant le rôle de ``flèches de
% archive sheet number: 269
\scanpage{270}
\mspage{220}
localisation'' ou ``d'iso. locaux'', satisfaisant les conditions suivantes

($\unsure{\alpha|a}$) $\unsure{\mathcal L|L}$ contient les id, est stable par composition, et par changement de base.
\inserted{On pose, pour $S$ dans $M$, $\mathcal L_S\subset M_{/S}$, sous-catégorie pleine des $S'_{/S}$ tels que $S'\longrightarrow S$ dans $\mathcal L$.}
On définit alors pour \deleted{\unsure{tout}} $S$ dans $M$
\[
\begin{aligned}
&W_S^{\mathcal L}\subset\mathrm{Fl}(M_{/S})\\
&\underset{\substack{\Vert\\(X\to Y)}}{u}\in W_S^{\mathcal L}
\Longleftrightarrow\quad\forall\, S'\in\mathcal L_S\,;\ \deleted{\unsure{\to}}\,
(u_{S'}:X\mathbin{\times_S} S'\unsure{-|\to}Y\mathbin{\times_S}S')\\
&\hspace{15em}\text{est dans }W.
\end{aligned}
\]
C'est \deleted{\unsure{donc}} une variante du localiseur $W_S^{\unsure{u|v}}$, où maintenant on ne prend comme chgts de base que les flèches \deleted{qui sont ``local} ``$\mathcal L$-localisantes'', i.e $\in\mathcal L$. On a donc
\[
W_S^u\subset W_S^{\mathcal L}\subset W_S
\]
en tant que parties de $\mathrm{Fl}(M_{/S})$, et de $=$ dans toute sous-catégorie pleine $\unsure{F|\mathcal F}$ de $M_{/S}$. Dans le cas de $\underline{\mathrm{Fib}}_W S$, moyennant l'axiome \textcircled{b} ci-dessus, ($W_S^u=W_S$ dans $\underline{\mathrm{Fib}}_W S$), les trois localiseurs coïncident.

On pose
\[
\deleted{\unsure{\mathrm{Ho}_W(S)=}}\mathrm{Ho}^{\mathcal L}_W(S)
\inserted{(\text{ou }\mathrm{Ho}_{W}(S))}
\overset{\text{déf}}{=}(M_{/S})(W_S^{\mathcal L})^{-1},
\]
et on a donné dans loc. cit (partie II du th. \unsure{salade}) \deleted{quelques} \inserted{un} résultat\deleted{s} intéressant\deleted{s} sur cet invariant, pour $S$ fixé. Mon propos
% archive sheet number: 270
\scanpage{271}
\mspage{221}
présent est de voir ce qui se passe pour $S$ variable. Visiblement, une notion de \inserted{$W$-}Type d'homotopie relatif sur un objet \unsure{quelc.} $S$ de $M$ n'est \unsure{satisfaisante|satisfaisant}, que si elle permet des images inverses par des flèches \uline{quelconques}
\[
f:S'\longrightarrow S
\]
dans $M$, qui doivent définir absolument
\[
f^{*\unsure{\mathcal L|L}}_{W}\ \text{ou}\ f^{\unsure{!|'}}_{W}:
\mathrm{Ho}_{W}^{\mathcal L}(S)\unsure{-|\longrightarrow}\mathrm{Ho}_{W}^{\mathcal L}(S').
\]
On est donc confronté à la $=$ perplexité que tantôt -- et ce n'est \uline{pas}, cette fois, l'introduction de $\underline{\mathrm{Fib}}_W S$ qui va sauver la situation, car
\[
(\underline{\mathrm{Fib}}_W S)
\bigl(\underset{\substack{\Vert\\W_S\\\Vert\\W_S^{\unsure{u|v}}}}{\underline{W_S^{\mathcal L}}}\bigr)^{-1}
\longrightarrow (M_{/S})(W_S^{\mathcal L})^{-1}
\]
n'a aucune envie, cette fois, d'être une équivalence de catégories, mais bien plutôt une \deleted{\unsure{f}} inclusion (\unsure{pl. fid}) des \inserted{$W$-}Types d'homotopie \unsure{rel|loc.} constants sur $S$, dans les $W$-types d'homotopie quelconques. Il faut visiblement un axiome qui \unsure{englobe|implique} \deleted{\unsure{aussi}} ``l'axiome \unsure{''|``}\deleted{\unsure{des}} de la fibre\deleted{\unsure{s}} homotopique''.
% archive sheet number: 271
\scanpage{272}
\mspage{222}
Dans le cas de $\mathrm{Cat}$, la classe $\unsure{\mathcal L|L}$ des iso. locaux (disons) permet de définir, comme ci-dessus, les localiseurs $\inserted{W_S^{\mathcal L}=}W_S^g$ dans les $\mathrm{Cat}_{/S}$, comme \unsure{formé|formée|formés} des $u:X\to Y$ sur $S$ qui sont dans $W$ et le restent par tt chgt de base $S'\to S$ qui est dans $\mathcal L$.
Mais il y a lieu de \uline{saturer} $\mathcal L$, en regardant \deleted{\unsure{pour tout}} la classe $\mathcal L^{\unsure{\wedge|1}}$ des flèches \inserted{$f:S'\to S$} de $\mathrm{Cat}$ telles que le chgmt de base
\[
f^*:\mathrm{Cat}_{/S}\unsure{-|\to}\mathrm{Cat}_{/S'}
\]
envoie $W_S^g$ dans $W_{S'}^g$. Ce $\mathcal L^{\wedge}$ est tautologiquement stable par composition, et contient les isos (mieux, contient $\mathcal L=\unsure{\text{ens.}}$ des iso. locaux) Il n'est pas clair, \unsure{sur} la définition donnée ici, qu'il soit stable par chgt de base -- on va \deleted{\unsure{changer}} \inserted{renforcer} la définition en conséquence, en exigeant, \deleted{\unsure{que}} \inserted{pour} que $f:S'\to S$ soit dans $\mathcal L^\wedge$, que pour tout $T$ sur $S$, considérons $f_T:T'=S'\times_S T\to T$,
$f_T^*:\mathrm{Cat}_{/T}\to\mathrm{Cat}_{/T'}$ \deleted{\unsure{on a|on ait}} soit compatible avec les localiseurs $W_{\unsure{T|T'}}^g$, $W_{T'}^g$. Alors \deleted{\unsure{pour f}} \inserted{$\mathcal L^\wedge$ est} stable \inserted{(tout comme $\mathcal L$)} par composition et par chgt de base, et \inserted{de plus il} contient $\mathcal L$. De plus, pour tt $S$, on a $W_S^{\mathcal L}=W_S^{\mathcal L^\wedge}$.
Ces \unsure{petits|petites} \unsure{sauts|sorties|soins} pouvaient se faire bien sûr pour \unsure{tt} modéliseur $(M,W)$. Mais la chose intéressante ici, c'est que $\mathcal L^\wedge$ contient l'\unsure{ens} \inserted{$\unsure{\mathrm{Liss}_W|\mathrm{Lisse}}$} des flèches $W$-\uline{lisses}. (J'ignore si elle est égale.) On \deleted{pourrait donc} \inserted{aurait donc} pu
% archive sheet number: 272
\scanpage{273}
\mspage{223}
aussi bien partir de l'\unsure{ens} de flèches $\mathrm{Liss}_W$, pour définir les localiseurs $W_S^g$ comme
\[
W_S^g=W_S^{\unsure{\mathrm{Lis}|\mathrm{Liss}}_W}\quad\text{au lieu de}\quad
W_S^g=W_S^{\unsure{\mathrm{Isoloc}|\mathrm{Iso\,loc}}},
\]
ce qui est équivalent.

D'autre part, en travaillant avec les \unsure{localisés} $\unsure{\mathrm{HOT}|\mathrm{Hot}}_W^0(S)=(\mathrm{Cat}_{/S})(W_S^g)^{-1}$, le rôle de $\underline{\mathrm{Fib}}_W S$ pour la théorie de $\mathrm{HOT}^{\unsure{\mathrm{lc}|\mathrm{li}}}_W(S)$ est joué ici par $\underline{\mathrm{Prop}}_W S$. On voit donc que dans la théorie des \inserted{$W$-}Types d'homotopie relatifs non loc. constants, notés $\mathrm{HOT}_W^0$, les ens. de flèches $\mathrm{Liss}_W$ et $\mathrm{Prop}_W$ jouent un rôle en quelque sorte complémentaire : les flèches \inserted{$f:S'\to S$} de $\mathrm{Liss}_W$ sont (pratiquement) celles qui permettent directement de définir
\[
f_W^{*0}:\mathrm{HOT}_W^0(S)\unsure{-|\longrightarrow}\mathrm{HOT}_W^0(S')
\]
comme induit par $f^*:(\mathrm{Cat}_{/S},W_S^g)\longrightarrow(\mathrm{Cat}_{/S'},W_{S'}^g)$,
tandis que les flèches de $\mathrm{Liss}_W$ permettent de définir \inserted{pour tt $S$ dans $\mathrm{Cat}$} la sous-catégorie pleine
\[
\underline{\mathrm{Prop}}_W S\subset\mathrm{Cat}_{/S}
\]
donnant lieu à une équivalence de \unsure{catégorie|catégories}
\[
\left\{\begin{aligned}
(\underline{\mathrm{Prop}}_W S)(W_S^{\unsure{u|v}})^{-1}&\xrightarrow{\approx}
\underset{\displaystyle =\,\mathrm{HOT}_W^0(S)}{\underline{(\mathrm{Cat}_{/S})(W_S^g)^{-1}}}\\
\text{car}\qquad &\\[-5pt]
W_S^u=W_S^g\ \text{sur } \underline{\mathrm{Prop}}_W S&
\end{aligned}\right.
\]
(ce qui visiblement est le substitut \unsure{annoncé} des ``axiomes de la \deleted{\unsure{co}} fibre $W$-homotopique''). Cette
% archive sheet number: 273
\scanpage{274}
\mspage{224}
équivalence, qui permet de décrire $\mathrm{Hot}_W^0(S)$ en termes de $\underline{\mathrm{Prop}}_W S\unsure{\subset|\subseteq}\mathrm{Cat}_{/S}$ et d'un localiseur $W_S^{\unsure{u|v}}$ \uline{stable par} \uuline{tout} \uline{chgt. de base}, permet d'autre part de définir la loi contravariante attendue pour $\unsure{\mathrm{HOT}|\mathrm{Hot}}_W^0(S)$.

D'autre part, quand on échange les rôles de $\deleted{\underline{\mathrm{Prop}}}\underline{\mathrm{Liss}}_W$ et $\mathrm{Prop}_W$, on trouve la théorie complémentaire des $\mathrm{HOT}_W(S)$,
\[
\left\{\begin{aligned}
\mathrm{HOT}_W(S)&\overset{\text{déf}}{=}(\mathrm{Cat}_{/S})(W_S^d)^{-1}
\xleftarrow{\approx}(\underline{\mathrm{Liss}}_W S)(W_S^u)^{-1}\\
\text{sur }\underline{\mathrm{Liss}}_W S,\quad &W_S^d=W_{\unsure{S|J}}^u,
\end{aligned}\right.
\]
la description de $\mathrm{HOT}_W(S)$ en termes de $\underline{\mathrm{Liss}}_W S$ et de $W_S^u$ permet d'en \deleted{établir} \inserted{définir} la variation contravariante en $S$, pour \uline{tte} flèche $f:S'\longrightarrow S$ dans $\mathrm{Cat}$. Tandis que si on veut pouvoir passer aux catégories de \unsure{fractions|foncteurs} sur $\mathrm{Cat}_{/S}$ et $\mathrm{Cat}_{/S'}$ tout entiers \inserted{(via $W_S^d$, $W_{S'}^d$)}, il faut se limiter aux chgts de base $W$-propres.

Du coup, me vient la perplexité si cette complémentarité, voire \inserted{cette} dualité, entre deux ensembles de flèches, les ``lisses'' et les ``propres'', ne serait pas un fait général qui se retrouverait tout au moins dans les catégories de modèles \unsure{proches|parties|pleines} de la \unsure{description} de $\mathrm{HOT}$ ou d'un $\mathrm{Hot}_W$.
% archive sheet number: 274
\scanpage{275}
\mspage{225}
Il se pourrait bien que dans $\Delta^\wedge$ on ait une dualité parfaite entre ces deux espèces de notions (il faudra que j'examine la chose), peut-être $=$ dans toute catégorie $A^\wedge$ (mais dans $\Delta^\wedge$ il y a un \unsure{autom.} involutif $K\mapsto K^{\unsure{\circ|0}}$, correspondant au $X\mapsto X^\circ$ dans $\mathrm{Cat}$ \ldots). \deleted{\unsure{Après}} Et après tout, même dans le contexte des espaces topologiques les notions de propreté d'une part, d'asphéricité locale d'autre part, pour une flèche $X\unsure{-|\to}Y$, sont bien définies, et pourraient bien s'ajuster comme le font les notions correspondantes dans $\mathrm{Cat}$, et (qui sait) donner lieu à deux versions complémentaires, ou ``duales'' l'une de l'autre, d'une notion de type d'homotopie relatif.

N'en ayant pas le cœur net, je suis un peu indécis comment présenter une description axiomatique de la théorie des $\mathrm{Ho}^{\unsure{\mathcal L|L}}_W(S)$, pour une \deleted{\unsure{local}} catégorie de modèles générale $(M,W)$. Dois-je me donner d'emblée
\[
\mathrm{Liss},\mathrm{Prop}\subset\mathrm{Fl}(M),
\]
\deleted{\unsure{avec}} stables par composition, par changement de base, contenant les iso, et tels que (peut-être)
\[
\mathrm{Liss}\cap\mathrm{Prop}\subset\mathrm{Fib}_W,
\]
% archive sheet number: 275
\scanpage{276}
\mspage{226}
mais surtout ils permettent de définir, pour \unsure{$S$|$C$} de $M$
\[
\left\{\begin{aligned}
W_S^g&\overset{\text{déf}}{=}W_S^{\mathrm{Liss}}\subset\mathrm{Fl}(M_{/S})\\
W_S^d&=W_S^{\mathrm{Prop}}\subset\mathrm{Fl}(M_{/S})
\end{aligned}\right.
\]
avec les propriétés suivantes :

\deleted{\textcircled{\unsure{b}}}\textcircled{a} Si $f:S'\to S$ est dans $\mathrm{Prop}$ \inserted{(resp. $\mathrm{Liss}$)}, alors
\[
\begin{aligned}
f^*:M_{/S}\longrightarrow M_{/S'}\quad &\text{transforme }W_S^d\text{ en }W_{\unsure{S|S'}}^d\\
&(\text{resp. } W_S^g\text{ en }W_{\unsure{S|S'}}^g)
\end{aligned}
\]
ce qui est une pure tautologie, et ne concerne pas encore les relations de $\mathrm{Prop}$ avec $\mathrm{Liss}$, mais par contre :
\[
\text{\textcircled{b}}\quad
\left\{\begin{aligned}
(\underline{\mathrm{Liss}}\,S)(W_S^{\unsure{u|v}})^{-1}&\xrightarrow{\approx}(M_{/S})(W_S^d)^{-1}
\quad\text{une équivalence}\\
\text{et}\qquad W_S^u&=W_S^d\quad\text{sur }\underline{\mathrm{Liss}}\,S
\end{aligned}\right.
\]
et dualement
\[
\left\{\begin{aligned}
(\underline{\mathrm{Prop}}\,S)(W_S^u)^{-1}&\xrightarrow{\approx}(M_{/S})(W_S^g)^{-1}
\quad\text{une \unsure{équivce|équiv\textsuperscript{ce}}}\\
\text{et}\qquad W_S^u&=W_S^g\quad\text{sur }\underline{\mathrm{Prop}}\,S.
\end{aligned}\right.
\]
De plus, il y a lieu d'énoncer \inserted{\textcircled{c}} \deleted{\unsure{une|une fois}} \inserted{deux} propriétés d'adjonction \inserted{(duales l'une de l'autre)}, pour lesquelles je renvoie à XII, théorème \unsure{salade} 2, p. 108-113. Il se pose ici la question subsidiaire de simplifier les conditions de validité de la formule d'adjonction, un peu hétéroclites dans loc. cit.

Il se pose \deleted{\unsure{en troi}} d'autre part la question de faire la théorie des \deleted{la} lois covariantes $f_!$ et $f_{\unsure{*|\star}}$ pour $\mathrm{Ho}^{\unsure{\mathrm{lc}|\mathrm{li}}}_W(S)$ de $=$ pour $\mathrm{Ho}^{\mathrm{Prop}}_W(S)$, $\mathrm{Ho}^{\mathrm{Liss}}_W(S)$.
% archive sheet number: 276
\scanpage{277}
\mspage{227}
J'ai bon espoir que \deleted{\unsure{l'on}} le foncteur $f_!$ peut être décrit, et son adjonction avec $f^*$ établie, à aussi bon compte dans ce contexte général que dans $\mathrm{Cat}$. Mais je n'ai pas envie de m'y arrêter à présent.

6. \inserted{\uline{Cofibre \inserted{$W$-}homotopique dans $\mathrm{Cat}$}}
Par contre, j'ai envie de regarder ici, \deleted{\unsure{sous}} sous quelles conditions sur le localiseur fondamental $W$, il existe dans $(\mathrm{Cat},W)$ une théorie duale de celle de la fibre $W$-homotopique.

La question-clef est, bien sûr, celle de \deleted{\unsure{exhiber}} trouver une bonne \unsure{prise sur les|provision de} $W$-cofibrations. Or je ne connais jusqu'à présent \uline{que} les immersions ouvertes de Dwyer, les immersions fermées de Dwyer, et ce qui s'en déduit par composition etc. Le travail de Thomason peut motiver l'espoir que \unsure{la connaissance|la considération|les conditions} de ces immersions-là, qui semblent ici jouer un rôle similaire \inserted{(dual)} à celui des morphismes $W$-parfaits, suffira pour faire une bonne théorie des \deleted{\unsure{la}} ``cofibres $W$-homotopiques''.

Je suppose dans \deleted{toute} cette section que $W$ satisfait $W(7\,\mathrm{bis})$ (axiome des carrés cocartésiens), pour avoir des bonnes propriétés des imm. ouvertes de Dwyer$^{*}$, et plus \unsure{gén\textsuperscript{t}|gén.} pour les \unsure{$W$-}cofibrations.
% margin note written sideways
\margin{a) Ce sont des $W$-cofibrations (résulte de b) et c))}
% margin note written sideways
\margin{b) Soit
\diagram{\begin{tikzcd}[ampersand replacement=\&,column sep=small,row sep=small]
\unsure{U|U'} \arrow[r,hook,"i"] \arrow[d,"p"'] \& \unsure{X|\alpha} \arrow[d,"q"]\\
\unsure{U'|W} \arrow[r,"i'"] \& \unsure{X'|X}
\end{tikzcd}}
avec $i$ imm. ouv. de Dw., alors
\[
\begin{aligned}
i\in W&\Longleftrightarrow i'\in W\\
p\in W&\Longleftrightarrow q\in W
\end{aligned}
\] (J'ai donc besoin de $W(7\,\mathrm{bis})$ \uuline{et} $\unsure{\supset|\Rightarrow}$ de $W(7)$ \inserted{pour les deuxièmes implications})}
% margin note written sideways
\margin{c) Les imm. ouvertes de Dw. sont stables par \unsure{cochangement} de base et par composition, contiennent les iso}
% archive sheet number: 277
\scanpage{278}
\mspage{228}
Il se pose donc la question de la factorisation duale d'une flèche \inserted{$f:X\to Y$} de $\mathrm{Cat}$ en
\diagram{\begin{tikzcd}[ampersand replacement=\&,column sep=large]
X \arrow[r,hook,"W\text{-cofibration}","i_f"'] \& \widetilde{\unsure{Y|X}} \arrow[r,"\in W","p_f"'] \& Y.
\end{tikzcd}}
\deleted{\unsure{Or Thomason dans}} \deleted{\unsure{Thom.}} Les résultats de Thomason assurent que ça existe, avec $i$ une imm. ouverte de Dwyer, et $p$ une $W$-fibration -- du moins dans le cas $W=W_\infty$. Je vais essayer de le voir directement, en calquant \inserted{dualement} la construction \deleted{\unsure{duale de celle}} de Cartan-Serre.

Il s'agit donc d'abord de factoriser
\diagram{\begin{tikzcd}[ampersand replacement=\&]
X\amalg X \arrow[r,"\mathrm{codiag}_X"] \& X,
\end{tikzcd}}
et pour ceci je choisis un intervalle fixe
\diagram{\begin{tikzcd}[ampersand replacement=\&,column sep=small]
I=\unsure{\Phi|\mathcal I}_n=\bigl(0 \arrow[r] \& 1 \& 2 \arrow[l] \& \cdots \& {} \arrow[r] \& 2n-1 \& 2n\bigr) \arrow[l]
\end{tikzcd}}
avec \qquad $n\geq 2$

de sorte que l'inclusion
\diagram{\begin{tikzcd}[ampersand replacement=\&,column sep=large]
\varepsilon=\begin{gathered}\text{catégorie}\\\text{discrète}\\\{0,1\}\end{gathered}
\arrow[r,"\partial","{\substack{\partial(0)=0\\\unsure{\pi|\partial}(1)=2n}}"'] \& I
\end{tikzcd}}
soit une immersion ouverte de Dwyer. Cette propriété est stable par chgt de base $\mathrm{Cat}$-fibrant, en particulier en prenant un produit
\diagram{\begin{tikzcd}[ampersand replacement=\&]
\underset{\simeq\,X\amalg X}{\underline{X\times\varepsilon}}\arrow[r,"\partial_X"] \& X\times I
\end{tikzcd}}
où
% archive sheet number: 278
\scanpage{279}
\mspage{229}
d'autre par\unsure{t}
\diagram{\begin{tikzcd}[ampersand replacement=\&]
X\times I \arrow[r,"\mathrm{pr}_1"] \& X
\end{tikzcd}}
est $W_{\unsure{\infty|w}}$-parfait, d'où la factorisation cherchée
\diagram{\begin{tikzcd}[ampersand replacement=\&]
X\amalg X \arrow[r,"i_X"] \& X\times I \arrow[r,"p_X=\mathrm{pr}_1^X"] \& X
\end{tikzcd}}
-- il n'y a pas plus simple ! Comment l'utiliser pour construire la factorisation d'une flèche $f:\unsure{\alpha|X}\to Y$ quelconque, en dualisant les constructions relatives aux espaces de chemins ? On définit $Z(f)$ par le carré cocartésien dans $\mathrm{Cat}$
\diagram{\begin{tikzcd}[ampersand replacement=\&,column sep=large]
X \arrow[r,dotted,"{\deleted{Z(f)}}","{\mathrm{inj}_{\unsure{1|2}}^X}"'] \& X\times I \arrow[r,"f'"] \& Z(f) \arrow[d,bend left=42,dashed,"{\unsure{P|p}_f}"]\\
\& X \arrow[u,hook,"{\mathrm{inj}_0^X}"'] \arrow[r,"f"] \& Y \arrow[u,hook,"{\unsure{w|W}}"']
\end{tikzcd},}
on a $Z(f)\xrightarrow{\alpha}Y$ induit par $\mathrm{pr}_1:X\times I\longrightarrow X$,
\margin{\figure{\begin{tikzpicture}[x=1ex,y=1ex]
\draw (0,2.9) -- (.16,2.75) -- (.20,2.55);
\draw (.56,2.67) .. controls (.95,2.69) and (.74,2.45) .. (.91,2.33) -- (.85,2.23)
.. controls (.98,1.87) and (.80,1.72) .. (.60,1.67) -- (.32,.66) -- (1.22,1.05);
\end{tikzpicture}}}
et la factorisation cherchée de $f$ est
\diagram{\begin{tikzcd}[ampersand replacement=\&,column sep=large]
X \arrow[r,"{i_f=f'\circ\mathrm{inj}_1^X}"] \& Z(f)\arrow[r,"P_f"] \& Y.
\end{tikzcd}}
Je dis que
\[
\left\{\begin{aligned}
i_f&\text{ est imm. ouverte de Dwyer, donc une}\\[-2pt]
&\qquad W\text{-cofibration}\\
p_f&\ \deleted{\unsure{\text{est }W\text{-asphérique}|\text{est }W\text{-parfait}}}\in W
\end{aligned}\right.
\]
% archive sheet number: 273
\scanpage{280}
\mspage{230}
On le voit sur le diagramme plus complet
\diagram{\begin{tikzcd}[ampersand replacement=\&,column sep=large,row sep=large]
\& X\times I \arrow[r,"f'"] \arrow[dd,dashed,bend left=48,"{\mathrm{pr}_1}" description] \& Z(f)\arrow[dd,bend left=40,dashed,"p_f"]\\
X\arrow[ur,dashed,"{\mathrm{inj}_1^X}"]\arrow[r,hook,"{\unsure{\mathrm{inj}|i}_1^X}"]
\& X\amalg X \arrow[u,hook,"i_X"'] \arrow[r,"{f\amalg\mathrm{id}_X}"]
\arrow[ur,phantom,"{\substack{\text{cocart et}\\\text{cart}}}" description]
\& Y\amalg X \arrow[u,hook,"j"']\\
\& X \arrow[u,hook,"{\unsure{i|\mathrm{inj}}_0^X}"'] \arrow[r,"f"]
\arrow[ur,phantom,"{\substack{\text{cocart et}\\\text{cart.}}}" description]
\& Y \arrow[u,hook]
\end{tikzcd},}
\[
i_f=f'\circ\mathrm{inj}_1^X=f'i_Xi_1^X
=j\,\underset{\substack{\text{inclusion can.}\\\text{de }X\text{ dans}\\Y\amalg X}}{\underline{(f\amalg\mathrm{id}_X)i_1^X}}
\]
est imm. ouverte de Dwyer, comme composée de $X\longrightarrow Y\amalg X$ et $j:Y\amalg X\longrightarrow Z(f)$ qui le sont ($j$ \inserted{obtenu} par cochangement de base à partir de $i_X$, qui est imm. ouverte de Dwyer).

\deleted{$\unsure{X\times I}$} D'autre part $X\xrightarrow{\mathrm{inj}_0^X}X\times I$ est imm. ouverte de Dwyer et dans $W$, donc couniversellement dans $W$, donc $Y\longrightarrow Z(f)$ est elle aussi dans $W$, donc aussi la rétraction $p_f$.

D'autre part, la proposition duale de
\[
\underline{\mathrm{Fib}}_W S\cap W_S=W^{\unsure{u|v}}_{S\,\unsure{\mathrm{univ.}|\mathrm{uni.}}}
\]
est valable pour les $W$-cofibrations :
% archive sheet number: 280
\scanpage{281}
\mspage{231}
% slanted
\margin{N’a rien à voir avec Cat -- c’est une chose \unsure{générale} quand il y a \unsure{assez} de $W$-cofibrations}
\uline{Proposition 16} Sur \deleted{\unsure{Cat|$ {}_{S\backslash}\mathrm{Cat}$}} la sous-catégorie pleine $\underline{\mathrm{Cofib}}_W(S)$ de $ {}_{S\backslash}\mathrm{Cat}$, le localiseur $ {}_S W$ est égal à $({}_S W)^{\unsure{\mathrm{univ.}|\mathrm{univ}}}$ ($W$-équivalences couniverselles), i.e. on a ceci : Soient
\[
S\longrightarrow X\qquad S\longrightarrow Y
\]
des $W$-cofibrations, et soit $f:X\longrightarrow Y$ un morphisme de $ {}_{S\backslash}\mathrm{Cat}$ tel que $f\in W$. Alors pour tout cochangement de base $S\longrightarrow S'$, la flèche correspondante
\[
f':X'=X\amalg_S S'\longrightarrow Y'=Y\amalg_S S'
\]
est dans $W$.

\uline{Dém} On factorise \deleted{\unsure{la|le}} $S\longrightarrow S'$ en
\[
S\xrightarrow{i}S\deleted{^{\unsure{0|\circ}}}'_0\xrightarrow{j}S',
\]
avec $i$ imm. ouv. de Dwyer (donc $W$-cofibration), et $j\in W$. On a donc le diagramme à carrés cocartésiens
\diagram{\begin{tikzcd}[ampersand replacement=\&]
S \arrow[r] \arrow[d,hook,"i"'] \& X \arrow[r,"f"] \arrow[d,hook] \& Y \arrow[d,hook] \\
S'_0 \arrow[r] \arrow[d,"{\in W}"'] \& X'_0 \arrow[r,dash,"{f'_0}","w"'] \arrow[d,"w"] \& Y'_0 \arrow[d,"w"] \\
S' \arrow[r] \& \unsure{x|X}' \arrow[r,"{f'}"] \& Y'
\end{tikzcd}}
Comme $i$ est une $W$-cofibration, \inserted{et $f\in W$,} on a $f'_0\in W$. D’autre part, comme $S\longrightarrow X$ donc $S'_0\longrightarrow X''_0$ est
% archive sheet number: 281
